\documentclass[11pt]{article}

\usepackage[margin=1in]{geometry}
\usepackage{amsmath,amssymb,amsthm}
\usepackage{mathtools}
\usepackage{stmaryrd}
\usepackage[T1]{fontenc}
\usepackage{xcolor}
\usepackage{hyperref}
\hypersetup{colorlinks=true,linkcolor=black,urlcolor=teal,citecolor=black}

\theoremstyle{plain}
\newtheorem{lemma}{Lemma}
\newtheorem{proposition}{Proposition}
\newtheorem{theorem}{Theorem}
\theoremstyle{definition}
\newtheorem{definition}{Definition}
\theoremstyle{remark}
\newtheorem{remark}{Remark}

% --- rho-calculus notation ---
\newcommand{\nil}{\mathbf{0}}
\newcommand{\for}[2]{\mathsf{for}(#1 \,{\leftarrow}\, #2)}
\newcommand{\out}[2]{#1\mathord{!}(#2)}
\newcommand{\drop}[1]{{*}#1}
\newcommand{\quo}[1]{@#1}
\newcommand{\dr}[1]{\,\overline{#1}\,}   % dereference / unquote (also written #x), set as a bar here
\newcommand{\subst}[2]{\{#1/#2\}}
\newcommand{\act}{\mathbin{*}}
\newcommand{\cmp}{\mathbin{\circ}}       % the name-pair "o"
\newcommand{\meas}{\mu}
\newcommand{\redto}{\longrightarrow}
\newcommand{\Proc}{\mathsf{Proc}}
\newcommand{\Names}{\mathsf{N}}
\newcommand{\FN}{\mathrm{FN}}
\newcommand{\mgu}{\mathrm{mgu}}
\newcommand{\lgg}{\mathrm{lgg}}
\newcommand{\unit}{\mathbf{1}}
\newcommand{\doi}[1]{\textsc{doi:}\,\texttt{#1}}

\title{\bf Fuzzy rho-calculus: a quantale-graded interaction,\\[2pt]
\large and its lift to ci\,GSLTs with binding}
\author{A working note}
\date{\today}

\begin{document}
\maketitle

\begin{abstract}
We present a concise axiomatization of a fuzzy rho-calculus and a minimal model: the comm rule
transmits not $Q$ but $(u\cmp x)\act Q$, where $u\cmp x$ is a quantale element measuring the match
between the two interaction surfaces and $a\act(-)$ is its action on terms. We compute the model in
full, show that the order axiom is a \emph{theorem} there (modulo one congruence side-condition
isolated as a lemma), and observe that the model commits to an \emph{accumulator} reading of $\cmp$,
with a dual \emph{attenuator} reading available. We then sketch the intended payoff: this is the flat-substitution,
asymmetric, single-surface instance of a general construction --- a quantale-graded interaction in a
continued--interactive GSLT --- and identify what each ingredient must become once continuations
carry binding, in both the asymmetric (``usual suspects'') and the symmetric (unification/fusion)
cases. The slogan that emerges: the comparison $\cmp$ is \emph{imposed} in flat rho but
\emph{generated} by unification in the symmetric lift, where $A$ is (a quotient of) the substitution
lattice and $\cmp$ is $\mgu$.
\end{abstract}

\section{Recap: rho and its fuzzification}

Recall the reflective syntax of the rho-calculus~\cite{rho}, with processes $P$ and names $x$
mutually defined, and the single structural reduction (comm):
\[
P,Q \;::=\; \nil \;\mid\; \for{y}{x}P \;\mid\; \out{x}{Q} \;\mid\; P\mid Q \;\mid\; \drop{x},
\qquad x,y \;::=\; \quo{P},
\]
\[
\for{y}{x}P \;\mid\; \out{x}{Q} \;\redto\; P\subst{\quo Q}{y}.
\]
Here $\quo{P}$ quotes a process to a name and $\drop{x}$ drops a name back to a process; we write
$\dr{x}$ for the dereference, or unquote (also written $\#x$), so $\dr{\quo P}=P$ and $\quo{\dr x}=x$.

\paragraph{Notational conventions.} We use $\Proc$ to denote the collection of all process terms generated by the grammar, and $\Names$ (or on occasion $\mathsf{@}\Proc$) to denote the collection of all names generated by the grammar.

\paragraph{The fuzzy data.} Fuzzification is purely semantic. One adds
\[
(-\act-)\colon A\times\Proc\to\Proc \quad(\text{action of grades on processes}),\qquad
(-\cmp-)\colon \Names\times\Names\to A \quad(\text{surface comparison}),
\]
and replaces strict channel equality in comm by a graded comm in which the transmitted term is
modulated by how well the two channels match:
\begin{equation}\label{eq:fcomm}
\for{y}{x}P \;\mid\; \out{u}{Q} \;\redto\; P\,\bigl\{\,\quo{\big((u\cmp x)\act Q\big)}\,\big/\,y\,\bigr\}.
\end{equation}
Read $u\cmp x$ as ``how closely $u$ and $x$ match'' and $(u\cmp x)\act Q$ as ``the part of (the
information in) $Q$ that gets through.'' Classical rho is the special case where $\cmp$ returns a
unit grade on a perfect match and $a\act(-)$ is then the identity.

\paragraph{The grade structure.} We demand that $A$ carries a quantale structure~\cite{mulvey,rosenthal}
$(A,\le,\otimes,\unit)$ --- a complete lattice with an associative product $\otimes$ that distributes
over arbitrary joins --- and a \emph{measure}
\[
\meas\colon \Proc \to A,
\]
subject to the
\begin{quote}\itshape
Order axiom.\quad $a_1\le a_2 \;\Longrightarrow\; \meas(a_1\act P)\le \meas(a_2\act P)$.
\end{quote}
A bigger match grade must read off as at least as much transmitted information.

\section{A minimal model}\label{sec:model}

The minimal model takes grades to be names themselves and realises every piece through
quote/par/dereference:
\begin{align}
A &= \Names, & x_1\cmp x_2 &= \quo{\big(\dr{x_1}\mid \dr{x_2}\big)},\notag\\
\meas(P) &= \quo P, & a\act P &= P\bigl\{\,(a\cmp z)/z \;:\; z\in\FN(P)\,\bigr\}.\label{eq:modeldefs}
\end{align}
The product is par-under-quote, the unit is $\quo{\nil}$ (since $P\mid\nil\equiv P$), the order is
sub-multiset inclusion of par-components, and joins are componentwise suprema of multiplicities.
Par is the \emph{product}, not the join; that distribution of $+$ over $\max$ is exactly what makes
this a quantale rather than a mere monoid.

\subsection{Unfolding the action}

Fix $u,x$ and $Q$ with $\FN(Q)=\{z_1,\dots,z_k\}$. Put $a=u\cmp x=\quo{(\dr u\mid \dr x)}$, so
$\dr a=\dr u\mid \dr x$. For each $z_j$,
\[
a\cmp z_j \;=\; \quo{\big(\dr a\mid \dr{z_j}\big)} \;=\; \quo{\big(\dr u\mid \dr x\mid \dr{z_j}\big)}.
\]
Hence the action is a uniform \emph{re-channelling}, never a deletion:
\begin{equation}\label{eq:unfold}
(u\cmp x)\act Q \;=\; Q\Bigl\{\,\quo{\big(\dr u\mid \dr x\mid \dr z\big)}\big/\,z \;:\; z\in\FN(Q)\Bigr\}.
\end{equation}
Every channel $z$ on which $Q$ acts is stamped with the \emph{match signature} $\dr u\mid \dr x$ of
the redex; the whole term structure of $Q$ survives intact.

\subsection{The two redexes}

With the example's orientation (input on $u_i$, output on $x_i$; commutativity of $\cmp$ makes the
orientation immaterial, since $\quo{(\dr u\mid \dr x)}\equiv\quo{(\dr x\mid \dr u)}$),
equation~\eqref{eq:unfold} gives
\begin{align*}
\for{y}{u_1}P \mid \out{x_1}{Q}
&\;\redto\; P\Bigl\{\,\quo{\big(Q\{\quo{(\dr{u_1}\mid \dr{x_1}\mid \dr z)}/z\}\big)}\big/\,y\Bigr\},\\[2pt]
\for{y}{u_2}P \mid \out{x_2}{Q}
&\;\redto\; P\Bigl\{\,\quo{\big(Q\{\quo{(\dr{u_2}\mid \dr{x_2}\mid \dr z)}/z\}\big)}\big/\,y\Bigr\}.
\end{align*}
The two residuals differ in exactly one place: the signature $\dr{u_1}\mid \dr{x_1}$ versus
$\dr{u_2}\mid \dr{x_2}$ stamped into every channel of $Q$.

\subsection{The order axiom is a theorem here}

\begin{lemma}[Congruence side-condition]\label{lem:cong}
Let the order $\le$ on $A=\Names$ be the least pre-congruence for the rho operators with
$\quo\nil$ as bottom (equivalently: sub-multiset inclusion at the top level, propagated through
prefixes $\mathsf{for}(\cdot)$, par, output, and quote). Then name-substitution into a fixed context
is monotone: $a\le a'$ implies $C[a]\le C[a']$ for every term context $C$.
\end{lemma}

\noindent The point of Lemma~\ref{lem:cong} is that a prefix $\for{y}{x}P$ is a single molecule;
sub-multiset comparison at the top level alone does \emph{not} push through it, and without
congruent propagation the monotonicity below leaks precisely at the binders.

\begin{proposition}[Grade monotonicity]\label{prop:mono}
In the model \eqref{eq:modeldefs} with the order of Lemma~\ref{lem:cong}, the order axiom holds.
Consequently, if $\dr{u_1}\mid \dr{x_1}\;\sqsubseteq\;\dr{u_2}\mid \dr{x_2}$ then
$\meas\big((u_1\cmp x_1)\act Q\big)\le \meas\big((u_2\cmp x_2)\act Q\big)$, i.e.\ the redex with the
larger match signature carries at least as much information into the residual.
\end{proposition}

\begin{proof}
$\cmp$ is $+$ on multisets, hence monotone in each argument; by Lemma~\ref{lem:cong} substitution of
a $\le$-larger grade into the fixed context $Q\{(-)/z\}$ yields a $\le$-larger term, so
$(u_1\cmp x_1)\act Q \sqsubseteq (u_2\cmp x_2)\act Q$. Applying $\meas=\quo{(-)}$, which is monotone
by construction, gives the claim. The order axiom is the same statement with $u\cmp x$ abstracted to
a free grade $a$.
\end{proof}

So in this model monotonicity is not an extra demand but a consequence of $\cmp$ being the (monotone)
quantale product and $\meas$ a monotone reading of it.

\subsection{What kind of ``fuzzy'' is this? Accumulator vs.\ attenuator}

\begin{remark}[The model commits to an accumulator]\label{rem:fork}
With $\cmp=\quo{(\dr u\mid \dr x)}$ the construction is an \emph{accumulator}, not an attenuator:
\begin{itemize}
\item it never deletes any of $Q$ --- by \eqref{eq:unfold} the action only renames channels, so all
term structure of $Q$ always survives;
\item $u\cmp x$ grows when \emph{either} side is large, regardless of overlap, so it measures combined
bulk, not similarity;
\item even a perfect match $u=x$ gives $u\cmp x=\quo{(\dr x\mid \dr x)}$, a doubling --- a
\emph{watermark}, not the identity. Classical rho is recovered only at $u=x=\quo\nil$.
\end{itemize}
If one wants the literal ``distance/attenuation'' reading --- perfect match transmits cleanly,
mismatch transmits less --- the natural swap is the multiset \emph{meet}
$u\cmp x=\quo{(\dr u\wedge \dr x)}$, the shared sub-process. Then $u\cmp x$ grows with genuine
overlap, and because meet can map distinct channels $\quo{Z_j}$ to a common stamped name, channels
can \emph{collapse}: attenuation as lost distinguishability rather than lost terms. Meet over
multisets is still a (frame-like) quantale and $\meas=\quo{(-)}$ stays monotone, so nothing breaks.
\end{remark}

\begin{remark}[A foretaste of the lift: $\mgu$ vs.\ $\lgg$]\label{rem:foretaste}
The fork in Remark~\ref{rem:fork} is not ad hoc. In the subsumption order on terms, the accumulator
(``keep everything from both surfaces'') is aligned with \emph{unification}~\cite{robinson} --- the most general common
\emph{instance}, $\mgu$ --- while the attenuator (``keep only what is shared'') is aligned with
\emph{anti-unification}~\cite{plotkin,reynolds} --- the least general common \emph{generalization}, $\lgg$. These are dual
operations in the same lattice. Section~\ref{sec:lift} takes this seriously: when continuations carry
binding and interaction is symmetric, the comparison $\cmp$ \emph{is} a (anti-)unifier and need not be
invented.
\end{remark}

\section{Lifting to ci\,GSLTs with binding}\label{sec:lift}

Rho is not binding-free --- the input prefix $\for{y}{x}P$ binds $y$ in $P$ --- but it has only that
one binding form, and in the model of \S\ref{sec:model} it is dispatched by ordinary capture-avoiding
substitution: the action \eqref{eq:modeldefs} acts on $\FN(P)$ as a flat set of free names, never
needing to be transported across $y$ in any non-trivial way. (Rho also lacks a $\nu$-style scoping
operator, so there is no scope to extrude.) In a continued--interactive GSLT, by contrast, a
\emph{language} is presented by a grammar (sorts/operations, several of which may bind), equations
(structural congruence), and rewrites (reduction)~\cite{nslogic,ntt,gslt}, with arbitrary declared
binders, and the continuation in an interaction redex \emph{carries binding structure}. The whole
graded apparatus must be reread inside the category of terms-with-binding --- nominal sets / a
presheaf category over name-contexts, or de~Bruijn families~\cite{gabbaypitts,pitts,fpt}. We list
what each of the three ingredients becomes; the flat rho action is then the degenerate case in which
the transport law below is trivial.

\subsection{(i) The action becomes a graded action respecting binders}

$a\act(-)$ is no longer a flat substitution but an $A$-graded, capture-avoiding endofunctor on the
term presheaf. The defining law is that it commutes with the binding (context-extension) functor:
writing $[y]\,P$ for the abstraction of $y$ in $P$ (whatever the binding form --- rho's input prefix,
$\lambda$, a $\nu$-restriction),
\[
a\act([y]\,P)\;=\;[y]\,(a'\act P),
\]
where $a'$ is $a$ transported across the binder (in nominal terms, $a$ taken with support away from
$y$, $\alpha$-renaming as needed). This transport law is the precise content of ``continuations have
binding structure.'' The order axiom lifts unchanged in statement: $\act$ is a lax action of the
quantale (functorial and monotone in $A$), in the sense of graded monads~\cite{katsumata}, and $\meas$
a lax monotone measure, whence \emph{fuzzy reduction is grade-monotone} --- now internal to the
presheaf category.

\subsection{(ii) The comparison $\cmp$: subject-matching vs.\ unification}

\paragraph{Asymmetric case (the usual suspects: $\lambda$~\cite{barendregt}, $\pi$~\cite{mpw}, rho).}
There is a privileged binding surface. The surfaces are a subject pair (the input channel $x$, the output
channel $u$); $\cmp$ compares subjects, $u\cmp x\in A$; and the action modulates the transmitted term
before a \emph{one-directional} graded substitution into the binder body, exactly as in
\eqref{eq:fcomm}. This is a near-verbatim lift: only the substitution becomes the presheaf's
capture-avoiding one, and the action must be transported under any binders occurring inside the
transmitted term.

\paragraph{Symmetric case (fusion/solos via unification).}
There is no privileged binder. Both surfaces present terms with (meta)variables, and the reaction is a
\emph{fusion}~\cite{parrowvictor,lanevevictor}: it computes a unifier $\sigma=\mgu(\text{surface}_1,\text{surface}_2)$ and applies $\sigma$
to \emph{all} continuations in scope, each application a graded action. The grade is read off the
unifier,
\[
\cmp(\text{surface}_1,\text{surface}_2)\;=\;\mathrm{grade}(\sigma)\in A,
\]
with natural choices ranging from the crudest two-element quantale
$\{\top\ \text{if unifiable},\ \bot\ \text{otherwise}\}$ (``fuzzy'' $=$ ``does it react at all'') to a
rich quantale built from the \emph{substitution lattice} --- substitutions ordered by instantiation,
$\mgu$ as its meet/join --- with $\mathrm{grade}(\sigma)$ a monotone invariant such as the number of
equations discharged, the term-size of $\sigma$, or $\meas$ of the substituted material.

\medskip\noindent This is the conceptual payoff. In flat rho the comparison $\cmp$ had to be
\emph{imposed} as extra structure (par-under-quote). Under symmetric interaction it is
\emph{generated}: the substitution lattice already is the requisite ordered structure, $\mgu$ already
is its meet/join, so $A$ may be taken to be (a quotient of) that lattice and $\cmp$ to be $\mgu$
itself. The intuition ``how far apart $u$ and $x$ are'' then becomes literal --- the
\emph{unification distance}: how much the $\mgu$ had to do to make the surfaces agree, or whether they
agree at all. Dually (Remark~\ref{rem:foretaste}), the attenuator/meet variant is realised by
anti-unification: $\cmp=\lgg$ keeps the shared skeleton of the two surfaces, the symmetric surface of the
$\quo{(\dr u\wedge \dr x)}$ swap.

\subsection{(iii) Side conditions: scope extrusion as a precongruence law}

Fusion extrudes scope: a fused name may sit under a restriction, and $\sigma$ must respect it (the
classical fusion side-condition --- do not fuse names trapped under conflicting binders; push the
restriction outward instead). The graded version adds a stability requirement that is the symmetric
analogue of Lemma~\ref{lem:cong}:
\[
\mathrm{grade}(\sigma)\ \text{must be stable under $\alpha$-conversion and scope extrusion,}
\]
i.e.\ computing the grade before and after pushing a $\nu$ outward must agree (up to the action's
transport law). Equivalently, the order on $A$ must be a pre-congruence for the binders, not merely
for par; otherwise grade-monotonicity leaks exactly at the binders --- the same failure
Lemma~\ref{lem:cong} forecloses in the flat, top-level case.

\section{Summary and open ends}

\noindent\textbf{Summary.} Fuzzy rho is the flat-substitution, asymmetric, single-surface, par-quote
instance of one construction: a quantale-graded interaction in a ci\,GSLT, with the grade obtained by
comparing interaction surfaces (subject-matching when asymmetric, (anti-)unification when symmetric),
acting capture-avoidingly on continuations, and grade-monotonicity guaranteed by a precongruent
order. The two warm-up redexes of \S\ref{sec:model} are the shadow this construction casts when
binding is trivial and the comparison is par.

\medskip\noindent\textbf{Open ends.}
\begin{enumerate}
\item Which monotone invariant of the $\mgu$ yields a $\meas$ for which grade-monotonicity says
something about \emph{genuinely transmitted} information, rather than mere substitution bulk?
\item Make the $\mgu/\lgg$ duality of Remark~\ref{rem:foretaste} a precise adjunction between the
accumulator and attenuator gradings, internal to the substitution lattice.
\item Reading the grade as a potential on the reduction graph (the energetics picture), is
grade-monotonicity along reduction a second-law-flavoured statement, and does the precongruence
condition correspond to a conservation law at the binders?
\end{enumerate}

\begin{thebibliography}{99}

\bibitem{rho}
L.~G. Meredith and M.~Radestock.
\newblock A reflective higher-order calculus.
\newblock \emph{Electronic Notes in Theoretical Computer Science}, 141(5):49--67, 2005.
\newblock \doi{10.1016/j.entcs.2005.05.016}.

\bibitem{nslogic}
L.~G. Meredith and M.~Radestock.
\newblock Namespace logic: A logic for a reflective higher-order calculus.
\newblock In \emph{Trustworthy Global Computing (TGC 2005)}, volume 3705 of \emph{LNCS}, pages 353--369. Springer, 2005.

\bibitem{ntt}
C.~Williams and M.~Stay.
\newblock Native type theory.
\newblock In \emph{Applied Category Theory (ACT 2021)}, volume 372 of \emph{EPTCS}, pages 116--132, 2022.
\newblock \doi{10.4204/EPTCS.372.9}.
\newblock Derives a type system from a language's presentation by passing to its presheaf topos (rho among the worked examples).

\bibitem{gslt}
L.~G. Meredith.
\newblock Generative spatial logic types (GSLT) and \textsc{MeTTaIL}: languages as types, equations, and rewrites.
\newblock Working notes / preprint, F1R3FLY.io.
\newblock \emph{[author to supply final reference]}.

\bibitem{mulvey}
C.~J. Mulvey.
\newblock \&.
\newblock \emph{Rendiconti del Circolo Matematico di Palermo}, Series II, Supplement 12:99--104, 1986.

\bibitem{rosenthal}
K.~I. Rosenthal.
\newblock \emph{Quantales and Their Applications}.
\newblock Volume 234 of Pitman Research Notes in Mathematics. Longman Scientific \& Technical, 1990.

\bibitem{robinson}
J.~A. Robinson.
\newblock A machine-oriented logic based on the resolution principle.
\newblock \emph{Journal of the ACM}, 12(1):23--41, 1965.

\bibitem{plotkin}
G.~D. Plotkin.
\newblock A note on inductive generalization.
\newblock \emph{Machine Intelligence}, 5:153--163, 1970.

\bibitem{reynolds}
J.~C. Reynolds.
\newblock Transformational systems and the algebraic structure of atomic formulas.
\newblock \emph{Machine Intelligence}, 5:135--151, 1970.

\bibitem{mpw}
R.~Milner, J.~Parrow, and D.~Walker.
\newblock A calculus of mobile processes, I and II.
\newblock \emph{Information and Computation}, 100(1):1--77, 1992.

\bibitem{barendregt}
H.~P. Barendregt.
\newblock \emph{The Lambda Calculus: Its Syntax and Semantics}.
\newblock Volume 103 of Studies in Logic. North-Holland, revised edition, 1984.

\bibitem{parrowvictor}
J.~Parrow and B.~Victor.
\newblock The fusion calculus: Expressiveness and symmetry in mobile processes.
\newblock In \emph{LICS 1998}, pages 176--185. IEEE, 1998.

\bibitem{lanevevictor}
C.~Laneve and B.~Victor.
\newblock Solos in concert.
\newblock \emph{Mathematical Structures in Computer Science}, 13(5):657--683, 2003.

\bibitem{gabbaypitts}
M.~J. Gabbay and A.~M. Pitts.
\newblock A new approach to abstract syntax with variable binding.
\newblock \emph{Formal Aspects of Computing}, 13(3--5):341--363, 2002.

\bibitem{pitts}
A.~M. Pitts.
\newblock \emph{Nominal Sets: Names and Symmetry in Computer Science}.
\newblock Volume 57 of Cambridge Tracts in Theoretical Computer Science. Cambridge University Press, 2013.

\bibitem{fpt}
M.~Fiore, G.~Plotkin, and D.~Turi.
\newblock Abstract syntax and variable binding.
\newblock In \emph{LICS 1999}, pages 193--202. IEEE, 1999.

\bibitem{katsumata}
S.~Katsumata.
\newblock Parametric effect monads and semantics of effect systems.
\newblock In \emph{POPL 2014}, pages 633--645. ACM, 2014.

\end{thebibliography}

\end{document}
